\section{Background \& Related Work}
\label{ch:background}

% ============================================================================
% === WRITING GUIDE: Chapter 2 (~8 pages; ~2,500 new words) =================
% ============================================================================
% JOB: prove the "two disconnected literatures" claim. Figurative-language
% NLP is detection-centric and never accessibility-oriented; E2R automates
% guidelines but has no figurative-language module. End on the measurement
% gap so that Ch 3 (what to build) and Ch 4 (how to judge it) are both
% pre-motivated.
% Sections 2.4 (Datasets) and 2.5 (Evaluation Methodology) are already
% written; new prose is needed for 2.1, 2.2 (new section), 2.3, and 2.6.
% All paper summaries live in wiki/thesis/papers/ (Obsidian vault).
% Drop Pukiene 2024 / Baytelman 2024 unless space allows a one-line mention.
% THEME HOOK: open with "two literatures, one missing guideline"; close 2.6
% with: the literature offers structure as a promising lever (Tian, Badran,
% Gao) and warns that no metric can adjudicate it (Lai & Nissim,
% Alva-Manchego, Scialom); the next two chapters take both seriously.
% ============================================================================

\subsection{Figurative Language}
% --- GUIDE (~1.5 pp): idioms vs metaphors as linguistic phenomena.
% Must land:
%  (1) Working definitions: idioms = conventionalised MWEs whose meaning is
%      non-compositional; metaphors = cross-domain mappings, identified at
%      token level via MIP/MIPVU \cite{mipvu}. Conventionality spectrum.
%  (2) Detection state of the art and its critique: \cite{Neidlein2020}
%      shows high F-scores mask conventionalised word-sense disambiguation,
%      not metaphor competence (frames the Ch 5 supervised comparison).
%  (3) LLM figurative competence is contested: \cite{Ichien2024} (GPT-4
%      emergent novel-metaphor interpretation) vs \cite{Dmitrijev2024}
%      (pessimistic, pre-GPT-4) vs \cite{Liu2022} (Fig-QA: above chance,
%      below human). Reconciliation: scale + task framing.
%  (4) Generation/adaptation side is thin: \cite{Lai2024} survey shows
%      evaluation is the bottleneck of figurative-language generation.

% TODO(jelle): prose.

\subsection{Easy-to-Read and Cognitive Accessibility}
% --- GUIDE (~1.5 pp): NEW SECTION, carries the FACILE positioning
% (Mari Carmen's deferred review item).
% Must land:
%  (1) E2R guidelines and their target populations; figurative language is
%      explicitly discouraged by E2R guidance, yet no guideline says how to
%      replace it: the missing module.
%  (2) The UPM per-guideline-module pattern: \cite{Surez-Figueroa2023}
%      (-mente adverbs, superlatives), \cite{Diab2024} (dialogue
%      conversion), \cite{Surez-Figueroa2024} (FACILE: assess + transform
%      architecture). Each module automates ONE guideline and validates with
%      the target population.
%  (3) Position this thesis as the figurative-language module of that
%      pattern, in English, with the Spanish port as future work (8.6).
%  (4) E2R vs generic simplification: distinct goals; rubric implications
%      (simplicity dimension is the E2R-critical one; pays off in Ch 7).

% TODO(jelle): prose.

\subsection{Large Language Models \& Agentic Systems}
% --- GUIDE (~1.5 pp):
% Must land:
%  (1) Zero/one/few-shot prompting; structured output; chain-of-thought.
%  (2) Decomposition precedents that make H2 a motivated hypothesis, not a
%      hunch: \cite{Tian2024} (theory-guided multi-step scaffolding beats
%      vanilla prompting on metaphor), \cite{Badran2025} (SemEval-2025
%      three-stage pipeline > monolithic on idioms), \cite{Gao2025} (PASS
%      paraphrase + self-check decomposition). Optionally \cite{Hayashi2025}
%      (context augmentation) as a related scaffolding pattern.
%  (3) The counterpoint: \cite{Jia2024} reaches SOTA metaphor detection with
%      supervised curriculum fine-tuning (supervised can still win; H1
%      cannot be assumed).
%  (4) Observability: multi-step systems expose intermediate state (traces),
%      monolithic prompting does not. One paragraph; examined qualitatively
%      in the Discussion (Section~\ref{sec:discussion-observability}).

% TODO(jelle): prose.

\subsection{Datasets}

\subsubsection{VU Amsterdam Metaphor Corpus (VUAMC)}

The VU Amsterdam Metaphor Corpus \cite{vuamsterdammetaphorcorpus} is a large, manually annotated corpus of English texts drawn from four genres: academic prose, news, fiction, and conversation. Metaphor annotation follows the MIP (Metaphor Identification Procedure) \cite{mipvu}: for each content word, annotators assess whether its contextual meaning contrasts with a more basic, concrete meaning from which the contextual meaning can be derived by comparison. Words meeting this criterion are labelled metaphorical.

The corpus is distributed as an XML file and contains 16{,}202 English sentences, each with per-token metaphoricity labels (binary: metaphorical / non-metaphorical). Crucially, the VUAMC provides \textit{token-level} annotations rather than span- or sentence-level ones, making it well-suited for evaluating fine-grained detection but challenging for span-based models, since a single conceptual metaphor may involve multiple individually annotated tokens.

A custom Python client was developed to parse the XML, extract sentence-token structures, and derive character-offset spans from the token annotations. All 16{,}202 sentences are ingested into the experiment database. The VUAMC does \textit{not} include gold literal paraphrases; replacement evaluation on this dataset therefore relies on detection metrics and human evaluation only.

\subsubsection{SemEval-2022 Task 2: Multilingual Idiomaticity Detection and Sentence Embedding}

SemEval-2022 Task 2 \cite{tayyar-madabushi-etal-2022-semeval} is a shared task on idiom detection and interpretation in English, Portuguese, and Galician. This thesis uses the English (EN) partition of Subtask A and Subtask B:
\begin{itemize}
    \item \textbf{Subtask A} provides sentences containing a pre-identified multi-word expression (MWE) and asks whether the MWE is used idiomatically in context (binary classification). Gold annotations include character-offset span boundaries for the MWE.
    \item \textbf{Subtask B} provides gold literal paraphrases of the full sentence for a subset of Subtask A examples, enabling evaluation of replacement quality.
\end{itemize}

The English partition contains 3{,}487 sentences after filtering and joining Tasks A and B. Gold paraphrases from Task B are stored alongside each example and used as references for BLEU and BERTScore evaluation. Unlike the VUAMC, SemEval 2022 Task 2 is evaluated at the sentence and span level, with no per-token labels.

\subsubsection{Other Figurative-Language Datasets}
% --- GUIDE (~0.25 pp): these were surveyed but NOT used in any experiment.
% Compress to one short prose paragraph ("other corpora considered and why
% the two above were selected"), not a "(Planned)" list: nothing is planned
% anymore. MAGPIE \cite{magpie} (idiom robustness), TroFi \cite{trofi} and
% MOH/MOH-X \cite{moh} (verb metaphor), MetaNet (conceptual grounding).
% Selection rationale: VUAMC = token-level gold + scale; SemEval = the only
% dataset with gold replacement paraphrases, which the thesis task requires.
\begin{itemize}
    \item MAGPIE \cite{magpie}: idiom interpretation robustness (qualitative / stress testing)
    \item TroFi \cite{trofi}: verb metaphor detection (focused evaluation)
    \item MOH / MOH-X \cite{moh}: verb metaphor detection (secondary validation)
    \item MetaNet: conceptual metaphor grounding (interpretation support)
\end{itemize}

\subsection{Evaluation Methodology}
\label{sec:eval-methodology}

Standard automatic metrics are commonly applied to generation and paraphrase tasks, but their suitability for figurative language is limited. This section outlines the key metrics used in this thesis and discusses their limitations in the context of figurative-to-literal adaptation.

\subsubsection{Automatic Metrics}

\paragraph{BLEU.} BLEU measures n-gram overlap between a system output and one or more reference texts~\cite{papineni2002bleu}. It is precision-oriented and applies a brevity penalty. For figurative language replacement, BLEU is severely misaligned: literal paraphrases of idiomatic expressions intentionally avoid lexical overlap with the gold reference, and multiple semantically equivalent paraphrases may share no surface n-grams. BLEU is retained here primarily as a baseline point of comparison, but its low values should not be interpreted as evidence of poor replacement quality.

\paragraph{SARI.} SARI decomposes quality into Add, Delete, and Keep operations relative to the input and references, and was originally designed for text simplification \cite{xu2016sari}. Figurative-to-literal transformation closely resembles simplification (replacing opaque expressions with transparent ones), making SARI conceptually well-motivated. It is nevertheless not used in this thesis: its reference dependence makes it inapplicable to the metaphor track (the VUAMC has no gold paraphrases), and on the idiom track it would be redundant with the human evaluation, which is the load-bearing instrument for replacement quality.

\paragraph{BERTScore.} BERTScore computes contextual token-embedding similarity between generated and reference texts using a pre-trained language model (RoBERTa in this work), enabling soft alignment rather than exact string matching \cite{zhang2020bertscore}. This makes it more robust to paraphrase variation than BLEU. However, high BERTScore does not guarantee correct interpretation: a plausible but semantically incorrect paraphrase may still achieve a high score. BERTScore is used as the primary automatic metric for replacement evaluation.

\subsubsection{Human Evaluation}

Automatic metrics alone are insufficient for assessing figurative-to-literal replacement quality. Recent work in sentence simplification has shown that most automatic metrics correlate only spuriously with human judgement on this class of task \cite{Alva-Manchego2021,Scialom2021}, motivating a primary reliance on human evaluation. The protocol adopted is \textit{Direct Assessment} \cite{Graham2013,Graham2017}, applied to the three-axis rubric (grammaticality, meaning preservation, simplicity) established for sentence simplification by Alva-Manchego et al.\ \cite{Alva-Manchego2020}. Human evaluation is reserved for the open-ended task of replacement, where no single ground truth exists; detection has gold annotations and is evaluated only with the automatic metrics above. The replacement evaluation covers idioms only (SemEval Task B); metaphors are evaluated through detection metrics on detect-then-replace runs (Section~\ref{sec:rq3}). The full survey methodology --- continuous-scale design, standardisation and aggregation, quality control, and three deliberate departures from the original DA specification --- is the subject of Chapter~\ref{ch:methodology}; the survey instantiation (system conditions, source-sentence pool, respondent recruitment) is part of the RQ3 experimental setup (Section~\ref{subsec:rq3-human-eval}).

\subsubsection{Multi-Layer Evaluation Framework}

The evaluation protocol is structured as four complementary layers, inspired by the framework proposed in \cite{Surez-Figueroa2024}:
\begin{enumerate}
    \item \textbf{Detection accuracy}: span-level precision, recall, and F1 (Section~\ref{sec:eval-framework});
    \item \textbf{Semantic similarity}: BERTScore and BLEU (supporting/automatic);
    \item \textbf{Interpretive correctness}: human evaluation of meaning fidelity (primary);
    \item \textbf{Conceptual grounding}: correctness of source--target metaphor mappings and idiom sense disambiguation (planned).
\end{enumerate}

\subsection{Synthesis: The Gap This Thesis Fills}
% --- GUIDE (~0.5 pp): NEW SECTION replacing the old "Literature Review"
% citation list (all of those papers are now cited in 2.1-2.3 above).
% Must land, in ~3 paragraphs:
%  (1) The coverage gap: the figurative-language cluster and the E2R cluster
%      do not intersect; no prior work treats figurative-language adaptation
%      as an E2R guideline module. That intersection is the contribution.
%  (2) The method bet: decomposition precedents (Tian, Badran, Gao) all
%      report wins on hosted closed-source backbones; whether they hold on
%      an open-weights deliverable is exactly what RQ2 tests.
%  (3) The measurement gap: no automatic metric is trusted for replacement
%      (Lai & Nissim; Alva-Manchego 2021; Scialom 2021), so the evaluation
%      must be human-grounded: Ch 4 builds that instrument.
% CLOSING THEME HOOK: "the literature offers structure as a promising lever
% and warns that no metric can adjudicate it; the next two chapters take
% both seriously."

% TODO(jelle): prose.

\newpage
